% Sources: related-work comparison/provenance, theory source cards and local primary texts.
\section{관련연구와 설계의 위치}
\label{sec:related}

기존 연구가 무엇을 제공하고, zkDPP가 그 위에서 무엇을
구체화하는지 살펴본다. 먼저 제품 주장의 근거를 검증하는 연구,
이어서 숨긴 이력을 조사하고 사용을 막는 연구를 비교한다.
마지막으로 같은 기능을 다른 기존 기술로 구성하는 경우를 검토한다.
이때 성질과 비용을 비교할 대상을 \emph{기준선}이라고 부른다.
논문에 어떤 절차가 없다고 해서 그 기술로 구현할 수 없다고
판단하지는 않는다.

\subsection{DPP 검증과 공급망 생애주기}

Udokwu와 Craß는 지속가능성 주장과 이를 뒷받침하는 비공개 자료를
분리하고 영지식증명으로 주장 조건을 검증하는 DPP 구성을
제안한다~\cite{udokwu-crass-2026}. Babel 등은 생산 단계의 검증 가능한
자격증명을 비공개 토큰의 이전과 결합하고, 사용한 지분을 소진하여
반복 사용을 방지하는 생애주기를 다룬다~\cite{babel-hicss-2025}.
이처럼 기존 DPP 연구도 여러 단계의 근거 연결과 사용을 다룬다.
본 연구는 여기에 사후 대응을 연결한다.
문제가 된 근거를 나중에 지정했을 때, 그것을 이어받은 기록 중
무엇이 남았고 무엇의 사용을 막아야 하는지 정하는 데 초점을 둔다. 회계 크레딧 배분과 디지털 소비 방지가 실제 재료의
구성이나 동일 물량의 외부 중복 등록까지 증명하는 것은 아니다.

Sahai 등의 공급망 시스템은 문서의 생성·소비 관계와 상태 전이 검증을
이용하여 오염 관련 이력 추적을 다룬다~\cite{sahai-supply-chain-2020}.
이는 응용과 전이 구조 측면에서 가까운 기준선이다. 다만 그 구성에서
입력 문서 목록은 공개되지만, zkDPP는 실제로 어떤 입력을 썼는지도
숨기려 한다. 문서 내용을 숨기는 것, 문서 사이의 연결을 숨기는 것,
소유자를 숨기는 것은 서로 다른 목표다. 따라서 각각을 나누어 비교한다.
zkDPP에서는 회계 상태를 한 번 소비하는 회계 증표·이전 증표와 외부 DPP의
지정 필드에 결합된 종단 제품 주장을 구분하고, 이를 감사의 탐색 대상 및
현재 차단 가능한 대상의 구분에 사용한다.

ZKVeil은 분산 식별자(decentralized identifier, DID)·자격증명과 비공개 제품·금액 조건의 증명을 양자 거래의
주문·지불·배송 상태에 결합한다~\cite{cao-zkveil-2026}. 거래의 제품
식별과 현재 세션 상태 검증은 본 연구에서도 고려해야 하는 기존 기능이다.
본 연구의 비교 대상은 그 위에 정의할 다단계 객체 이력의 복원과 전체
후속 미소비 집합의 대응 절차이며, 개별 세션의 정상 완료와 이를 구분한다.

\begin{table*}[t]
\centering
\small
\caption{Closest mechanisms and the scope of comparison}
\label{tab:related}
\begin{tabularx}{\textwidth}{@{}p{30mm}YY@{}}
\toprule
Baseline & Related mechanism & Distinction for comparison with zkDPP \\
\midrule
Udokwu--Craß~\cite{udokwu-crass-2026}
& Private-input sustainability predicates
& Accepted lifecycle graph and live-object response \\
Babel et al.~\cite{babel-hicss-2025}
& Linked credentials, private ownership transfer, fraction consumption
& Post hoc targets and current consumption checks \\
Sahai et al.~\cite{sahai-supply-chain-2020}
& Supply-chain consume/create rules and public input document references
& Hidden actual parent correspondence and authorized recovery \\
GGM~\cite{garman-accountable-privacy-2017}
& Proof-bound tracing data and expansion to descendant outputs
& Typed lifecycle, enforcement targets and completion \\
Azeroth~\cite{jeong-azeroth-2023}
& Transaction-bound encrypted audit information
& Explicit graph traversal and whole-target response \\
Lee et al.~\cite{lee-token-guided-2026}
& Flow-, direction- and epoch-scoped audit capability
& Key-domain-wide auditor access; separate containment procedure \\
CAP~\cite{espresso-cap-spec-2022}
& Policy-controlled viewing and freezing; consume-and-reissue freeze
& Full-reveal/threshold conditions; shared spend ID versus encrypted future ID \\
ERC-20R~\cite{wang-erc20r-2022}
& Downstream amount tracing and atomic freezing on public history
& Account liability versus full descendant sets; atomic execution versus spending races \\
ACTMeds~\cite{sitaula-actmeds-2026}
& Private recall-reason proofs; transitions restricted by batch status
& Batch-level restrictions versus all recovered live descendants \\
ZEXE family and composed baseline~\cite{bowe-zexe-2020,xiong-verizexe-2023}
& Generic private creation and consumption rules
& Same trust, key access and leakage; same application relations \\
\bottomrule
\end{tabularx}
\end{table*}

\subsection{비공개 이력의 감사와 사용 제한}

이력의 열람 권한과 후속 사용을 막는 권한은 서로 다른 기능이다.
추적 자료를 실제 거래에 연결하는 방식, 승인된 공개의 범위와
동결의 집행 대상을 함께 비교한다.

GGM은 비공개 지급의 추적 정보를 증명 관계에 결합하고, 추적 대상
입력이 소비되었을 때 후속 출력으로 추적을 이어가는 구성을
제시한다~\cite{garman-accountable-privacy-2017}. 실제 출력과 연결된
암호화된 추적 정보 또는 후속 추적 자체를 본 연구의 새로운 원시기술로
볼 수 없는 이유이다. Azeroth 역시 감사자가 거래의 감사 정보를
복호화할 수 있도록 하고 그 정보와 실제 거래의 결합을
다룬다~\cite{jeong-azeroth-2023}. 따라서 Azeroth와 비교할 때 감사자가 한 거래만
열 수 있다고 가정하지 않는다.

Lee 등의 Token-Guided Flow Tracing은 방향별 추적 토큰과 에포크
키를 이용하여 승인된 거래 흐름·방향·기간 안에서 감사 정보를
열도록 구성한다~\cite{lee-token-guided-2026}. 토큰 관리자와
감사자가 승인 범위를 넘어 정보를 얻으려고 협력하지 않는다고 가정하며, 수락된 노트의 소비 관계와
사용자별 연속 거래를 따라 추적한다. zkDPP는 유효한 감사 데이터
키를 확보한 신뢰된 감사자가 같은 키 영역의 기록을 복호화하는
모델이므로 이와 같은 암호학적 범위 제한을 제공하지 않는다.
감사의 공개 범위와 그 결과를 사용 차단으로 연결하는 절차는
서로 다른 비교 항목이다.

Derecho는 증명 동반 공개(proof-carrying disclosure)로 허용된 출처에 관한 증명을
거래의 실제 입력을 따라 재귀적으로 전파한다~\cite{beal-derecho-2023}.
이는 출처 조건의 증명과 거래 관계의 결합에 관한 선행기술이며,
기존에 만들어진 증명을 사후 정책 변경에 따라 취소하는 데 따르는
한계도 논의한다. zkDPP는 문제 시작점을 나중에 지정한 후 수락된
객체 그래프를 감사자가 복원하는 절차를 다룬다. 기존 제품 주장을 취소하거나
과거 증명을 소급하여 무효화하지 않고, 현재 미소비 객체의 상태를
이후 소비 시점에 검사한다. 이 차이는 비교하는 보장의 범위이지
Derecho의 일반적 표현력에 대한 제한이 아니다.

이력을 조사한 뒤 어떤 대상의 사용을 제한할 수 있는지는
각 구성의 동결·리콜 절차에서 정해진다.

CAP는 비공개 자산에 대해 정책별 입력·출력 열람과 동결 권한을
정의한다~\cite{espresso-cap-spec-2022}. 열람용 메모(viewer memo)와 거래의 결합,
소유자와 동결 권한자가 같은 소비 식별자를 계산할 수 있는 구조,
동결 권한자에 의한 레코드 소비·재발행은 zkDPP와 직접 비교해야 하는
요소이다. 열람 범위는 정책과 조건에 따르며, 동결은 하나 이상의
레코드를 동결 상태로 재발행하는 방식이다. zkDPP의 현재 식별자에
대한 상태 변경과 구체 전이는 다르지만, 비공개 거래의 열람과 사용
차단을 결합한다는 일반적 기능은 이미 CAP에 존재한다.

CAP의 동결 정책은 전체 레코드 정보의 열람을 요구하며,
열람·동결 정책은 설정된 거래량 임곗값의 적용을 받는다. 여러
입력을 한 동결 트랜잭션에서 소비·재발행할 수 있다는 점도 비교에
포함해야 한다. zkDPP의 미래 식별자 사전 암호화와 CAP의 소유자·동결
권한자 간 식별자 계산은 집행 정보를 얻는 서로 다른 구성이다.
전자는 출력 생성자의 비밀 지식과 이전 증표 전달을 요구하며,
후자보다 자동으로 적은 공개나 비용을 제공한다고 볼 수 없다.

ERC-20R은 공개 거래 이력에서 후속 금액을 추적하고 하나의 원자적
동결 트랜잭션으로 회수 대상 금액을 확보하는 절차를
제안한다~\cite{wang-erc20r-2022}. 이 방식은 공개 계정에서 어느 금액을 회수할지 정한다.
반면 zkDPP는 원료가 합쳐지거나 나뉜 뒤의 모든 후속 출력을
따라가 남은 기록의 집합을 정한다. 대상의 단위가 서로 다르다. 특히 여러
상태 변경을 확인하는 zkDPP 실행에는 소비 경합과 부분 완료가 있으므로,
ERC-20R의 원자성을 그대로 보장한다고 서술하지 않는다.

ACTMeds는 영지식증명으로 리콜 사유를 검증하고 의약품 배치가
리콜 상태가 된 뒤의 전이를 제한하는 설계를 기술한다~\cite{sitaula-actmeds-2026}.
공개된 설계 설명을 기준으로 한 비교이며, 본 연구에서 그 실행 코드의
적합성을 검증한 것은 아니다. 여기서 배치 상태에 의한 리콜과
감사로 복원한 전체 후속 회계 증표·이전 증표의 식별은 구분해야 한다.
Circularise의 공개 특허는 비공개 의존 관계와 정보 유효성 유지에 대해
만료·취소·기록 동결 등의 대안을 논의한다~\cite{circularise-wo2025056975-2025}.
그중 기록 동결은 기록의 추가 갱신을 막는 의미이며, 디지털 자산의
소비 차단과 동일하지 않다. 특허의 여러 대안을 하나의 구현이 모두
제공한다는 결과로 해석하지 않는다.

\subsection{구성 대안과 설계 선택}
\label{sec:composed-baseline}

ZEXE와 VeriZexe는 비공개 상태의 생성·소비 규칙을 표현하는 범용
기반을 제공한다~\cite{bowe-zexe-2020,xiong-verizexe-2023}.
zkDPP의 전이, 감사 암호문 결합, 상태 검사도 이러한 기반 위에서
구성할 수 있다. 이러한 대안과 비교하려면 감사자의 신뢰 범위,
키의 열람 범위, 공개 정보와 회계·소비·완료 규칙을 맞춰야 한다.
아래의 조합 기준선은 이 조건을 맞춘 분석용 구성이다.

Table~\ref{tab:related}는 이 기준으로 가장 가까운 기능을 요약한다.
각 행은 기능의 유무를 단정하는 목록이 아니라 집행 대상과 가정의
차이를 보여 준다. zkDPP가 제시하는 것은 생산·유통 객체의 실제 전이,
승인된 이력 복원, 미소비 집합의 현재 사용 차단을 함께 명세한 DPP
구성과 그 완료 조건의 분석이다.

조합 기준선 $C^\star$는 범용 비공개 생성·소비 관계에 본 논문과
동일한 회계 정책, 감사 키 접근 및 공개 의미를 부여한 분석용
구성이다. 그 관계에 실제 부모·출력 식별자의 암호화를 넣고,
원장이 타입별 소비 기록과 상태를 검사하며, 감사자가 같은 전체
대상과 완료 조건을 사용하면 \ref{sec:analysis}절의 논증은
$C^\star$에도 적용된다. 따라서 동일 기능의 기준선과 zkDPP 사이에
기능적 우위를 설정하지 않는다.

zkDPP는 이 연결을 세 가지 선택으로 구체화한다.
첫째, 회계 증표를 앞으로 사용할 소유자가 직접 생성하고,
조직 간 인계는 이전 증표로 분리한다. 송신자가 수신자의 소유권
비밀을 몰라도 각 기록의 생성자가 미래 소비 식별자를 계산해
감사 자료에 남길 수 있다. 대신 송신과 인수의 두 거래 및
이전 증표 개방값의 비밀 전달이 필요하다.

둘째, 감사 평문에 기록의 종류·위치와 미래 소비 식별자를 결합한다.
종류·위치는 과거 기록을 찾는 데, 미래 식별자는 소비 여부의 확인과
현재 상태 변경에 쓰인다. 부모 참조만 읽는 구성은 미소비 객체의
집행 식별자를 별도로 확보해야 한다. zkDPP는 기록마다 이를 남겨
직접 탐색을 지원하며, 그만큼 암호문 저장과 검증 비용을 부담한다.
전체 자료를 모아 조회표를 만드는 대안은 \ref{sec:limits}절에서 다룬다.

셋째, 정확한 추적과 전체 차단 완료를 따로 판정한다.
여러 트랜잭션의 부분 성공을 전체 성공으로 오인하지 않도록
대상 목록을 보존하고 공통 시점의 상태를 확인한다.
추가 조회와 결과 확인이 필요하며, 소비 경합이 있으면 완료 대신
미완료를 반환한다. 이 선택과 그 보장 조건은
$C^\star$를 포함해 같은 규칙을 따르는 구성에도 적용할 수 있다.
본 논문의 기여는 이 구체적인 구성과 조건의 명세·분석이며,
새 암호 원시기술이나 다른 구성보다 높은 효율을 제시하는 것은 아니다.
